\begingroup
\renewcommand{\arraystretch}{1.45}
\newcommand{\TransscendentArticle}[1]{\par\medskip\Needspace{5\baselineskip}\pdfbookmark[1]{Art. #1}{transscendentes-ii-art-#1}\centerline{[#1.]}\nobreak\smallskip}
\fancyhead[L]{\small THEORY OF THE NEW TRANSCENDENTAL FUNCTIONS. II.}

\Needspace{20\baselineskip}
\SourceLeaf{450}
\TransscendentArticle{5}
The following theorem is one of the principal theorems of this theory.
If one denotes
\[
\begin{gathered}\{e^{-\frac{P\alpha}{24}}-e^{-\frac{25P\alpha}{24}}
-e^{-\frac{49P\alpha}{24}}\\[7pt]\quad{}+e^{-\frac{121P\alpha}{24}}
+e^{-\frac{169P\alpha}{24}}+\mathinner{..}\}\mathbin{.}{}^{4}\!\surd\alpha
\quad\text{by }\varphi\alpha\end{gathered}
\]
then
\[
\varphi\alpha=\varphi\frac{1}{\alpha}
\]
\Needspace{8\baselineskip}
This function reaches its maximum for $\alpha=1$, where its value
\[
=0{,}7682255=\frac{1}{{}^{4}\!\surd2\mathbin{.}\surd1{,}19814\mathinner{..}}
={}^{4}\!\surd\tfrac12\mathbin{.}\surd\frac{\varpi}{\pi}
\]
furthermore,

\Needspace{8\baselineskip}
\[
(\varphi 2\lambda)^{24}=4(\varphi\lambda)^8(\varphi 4\lambda)^8
\{(\varphi\lambda)^8+(\varphi 4\lambda)^8\}
\]

\Needspace{8\baselineskip}
\[
\varphi\alpha=e^{-\frac{\alpha}{24}P}{}^{4}\!\surd\alpha\mathbin{.}
(1-e^{-aP})(1-e^{-2\alpha P})(1-e^{-3\alpha P})\ .\ .\ .
\]
or

\Needspace{10\baselineskip}
\[
[x]=\frac{\varphi\frac{-\log x}{P}\mathbin{.}{}^{24}\!\surd\frac1x}
{{}^{4}\!\surd\frac{-\log x}{P}}
\]
\Needspace{13\baselineskip}
Let
\[
\begin{gathered}(1+xy)(1+x^3y)(1+x^5y)\ .\ .\ .
\\[7pt]\quad{}(1+\tfrac{x}{y})(1+\tfrac{x^3}{y})(1+\tfrac{x^5}{y})\ .\ .\ .=F(x,y)\end{gathered}
\]
then

\Needspace{13\baselineskip}
\[
\begin{gathered}F\bigl[e^{-\frac12\alpha P},\quad e^{P\omega\surd\alpha}\bigr]
\mathbin{.}e^{-\frac{P}{24\alpha}}
\\[7pt]\quad{}=F\bigl[e^{-\frac12\frac{P}{\alpha}},\quad
e^{iP\omega\surd\frac1\alpha}\bigr]
\mathbin{.}e^{-\frac{P\alpha}{24}}\mathbin{.}e^{\frac12P\omega\omega},\end{gathered}
\]


\Needspace{20\baselineskip}
\TransscendentArticle{6}
The division into seven parts leads to the following equation
\[
\begin{gathered}
b=\Bigl[\frac{1-2x+2x^4-\mathinner{..}}{1+2x+2x^4+\mathinner{..}}\Bigr]^2\\[7pt]
A=\Bigl[\frac{1+2x^7+2x^{28}+\mathinner{..}}{1+2x+2x^4+\mathinner{..}}\Bigr]^2\\[7pt]
B=\Bigl[\frac{1-2x^7+2x^{28}+\mathinner{..}}{1+2x+2x^4+\mathinner{..}}\Bigr]^2
\end{gathered}
\]


\Needspace{18\baselineskip}
\SourceLeaf{451}
\[
\begin{gathered}
A^8-\tfrac47A^6+\frac{16-32bb}{49}A^5-\frac{30}{343}A^4
\\[7pt]\quad{}+\frac{32-64bb}{2401}A^3
\\[7pt]\quad{}-\frac{20+768bb-768b^4}{16807}A^2\\[7pt]
+\frac{48-2144bb+6144b^4-4096\mathbin{.}b^6}{823543}A
-\frac1{823543}=0
\end{gathered}
\]

\Needspace{18\baselineskip}
\[
\begin{gathered}
B^8-\tfrac47bbB^6+\frac{16b^3-32b}{49}B^5-\frac{30b^4}{343}B^4
\\[7pt]\quad{}+\frac{32b^5-64b^3}{2401}B^3
\\[7pt]\quad{}-\frac{20b^6+768b^4-768b^2}{16807}B^2\\[7pt]
+\frac{48b^7-2144b^5+6144b^3-4096b}{823543}B
-\frac1{823543}=0
\end{gathered}
\]
\Needspace{8\baselineskip}
If, in the first equation, one replaces $bb$ by $1-bb$ and $A$ by $-A$, the equation is unchanged.


\Needspace{20\baselineskip}
\TransscendentArticle{7}
For the division into three parts, one has
\[
\begin{array}{ccc}
a&b&\dfrac ba=t\\[5pt]
A&B&\dfrac BA=T
\end{array}
\]
\[
\begin{gathered}T^4+(12t-16t^3)T^3+6ttT^3\\[7pt]\quad{}-(16t-12t^3)T+t^4=0\end{gathered}
\]
or
\[
(T-t)^4=16(t-t^3)(T-T^3)
\]
\[
\frac Aa=\frac{3t-t^3+T(1-3tt)}{3[T(1-tt)+t(1-TT)]}
\]
\Needspace{13\baselineskip}
If one sets $T=\operatorname{tang}N$ and $t=\operatorname{tang}n$, then
\[
\begin{gathered}
3\sin(2N-2n)^2=4\sin(N+3n)\mathbin{.}\sin(3N+n)\\[7pt]
\sin(N-n)^4=\sin4n\mathbin{.}\sin4N
\end{gathered}
\]

\Needspace{10\baselineskip}
\[
\begin{gathered}\frac{A\cos n}{a\cos N}=\frac{B\sin n}{b\sin N}
\\[7pt]\quad{}=\frac{2\sin(N+3n)}{3\sin(2N-2n)}
=\frac{\sin(2N-2n)}{2\sin(3N+n)}\end{gathered}
\]

\clearpage
\endgroup
